\documentclass[10pt,a4paper]{amsart}
\usepackage[margin=19mm]{geometry}
\usepackage{amsmath,amssymb,mathtools}
\usepackage[scr=rsfs,cal=euler]{mathalfa}
\usepackage{xfrac}
\usepackage[unicode,hidelinks,bookmarks=false]{hyperref}
\newcommand{\ie}{i.e.\ }
\newcommand{\cf}{cf.\ }
\newcommand{\resp}{respectively\ }
\newcommand{\Subsection}{\subsection}
\providecommand{\nobreakdash}{\nobreak\hbox{-}}
\setlength{\emergencystretch}{2em}
\multlinegap0pt
\renewcommand{\P}{\mathbb{P}}
\newcommand{\D}{\mathbb{D}}
\newcommand{\ED}{\mathrm{ED}}
\newcommand{\W}{\mathbb{W}}
\newcommand{\F}{\mathcal{F}}
\newcommand{\essinf}{\mathrm{ess\,inf}}
\newcommand{\eE}{\mathcal{E}}
\newcommand{\G}{\mathcal{G}}
\newcommand{\M}{\mathbb{M}}
\newcommand{\R}{\mathbb{R}}
\newcommand{\N}{\mathbb{N}}
\newcommand{\Z}{\mathbb{Z}}
%\renewcounter{equation}[section]
\newcommand{\E}{\mathbb{E}}
\newcommand{\1}{\mathbf{1}}
\newcommand{\ust}{\mathbf{u}}
\newcommand{\bs}{\mathbf{s}}
\newcommand{\dd}{\mathrm{d}}
\newcommand{\bv}{\mathbf{v}}
\newcommand{\cC}{\mathcal{C}}
\newcommand{\tr}{\operatorname{tr}}

\newcommand{\coloredZ}{\mathcal{Z}_{G}^{\mathcal{C}}}

\def\nb{{\mathfrak{nb}}}
\def\onb{\overline{\mathfrak{nb}}}
\def\pa{{\mathfrak{pa}}}
\def\opa{\overline{\mathfrak{pa}}}
\def\ch{{\mathfrak{ch}}}
\def\och{\overline{\mathfrak{ch}}}
\def\an{{\mathfrak{an}}}
\def\oan{\overline{\mathfrak{an}}}
\def\de{{\mathfrak{de}}}
\def\ode{\overline{\mathfrak{de}}}
\def\nd{{\mathfrak{nd}}}

\newcommand{\nc}{\newcommand}
\nc{\ep}{\varepsilon}
\nc{\Zsp}{\mathcal{Z}}
\nc{\Lsp}{\mathcal{L}}
\nc{\Ssp}{\mathcal{S}}
\nc{\Msp}{\mathcal{M}}
\nc{\Pcone}{\mathcal{P}}
\nc{\BC}{BC}
\nc{\cBC}{DCBC}
\nc{\IG}{$I$}
\nc{\cIG}{$cI$}
\nc{\hh}{\mathfrak{h}}
\newcommand{\HI}[1]{\textcolor{blue}{#1}}
\newcommand{\Delete}[1]{\textcolor{yellow}{[To Be Deleted]#1}}
\newcommand{\mytransp}[1]{ #1^{\top}}

\newcommand{\bdiag}{\mathrm{BlockDiag}}
\newcommand{\tri}{\mathrm{BlockTri}}
\newcommand{\RefExt}{\mathrm{RefExt}}

\def\ess{{\mathrm{ess}}}
\def\supp{{\mathrm{supp}}}

\theoremstyle{plain}
\newtheorem{theorem}{Theorem}
\newtheorem{propo}[theorem]{Proposition}
\newtheorem{lemma}[theorem]{Lemma}
\theoremstyle{definition}
\newtheorem{definition}[theorem]{Definition}
\newtheorem{example}[theorem]{Example}
\theoremstyle{remark}
\newtheorem{remark}[theorem]{Remark}
\newtheorem{corol}[theorem]{Corollary}
\title[A new class of colored Gaussian graphical models with explicit normalizing constants]{A new class of colored Gaussian graphical models with explicit normalizing constants: the characterisation of RDAG-admitting colored graphs (Theorem 6.1)}
\author{Adam Chojecki, Piotr Graczyk, Hideyuki Ishi, Bartosz Ko{\l}odziejek}
\date{Source: arXiv:2601.16945v2 (CC BY 4.0)}
\begin{document}
\maketitle

\noindent\emph{Source and licence.} A new class of colored Gaussian graphical models with explicit normalizing constants. Version arXiv:2601.16945v2 (CC BY 4.0), distributed under the Creative Commons Attribution 4.0 International licence, http://creativecommons.org/licenses/by/4.0/, as recorded in the arXiv licence metadata for this exact version and shown on the abstract page https://arxiv.org/abs/2601.16945v2. Full licence notice: licenses/arxiv-2601.16945-CC-BY-4.0.txt.\par\smallskip

\noindent\textit{Editorial scope.} Counted unit: the article's Theorem 6.1 (source0.tex lines 2264--2270, source label \texttt{thm:RDAG}), the characterisation of the colored graphs that admit an acyclic orientation with compatible inherited coloring and no induced subgraph $i\to k\leftarrow j$ with $i$ and $j$ non-adjacent, namely exactly the symmetric CER graphs with no edge inside a vertex color class; it is delivered with the article's complete original proof (source0.tex lines 3490--3536, one \texttt{proof} environment of 47 lines carrying the optional argument ``Proof of Theorem~\ref{thm:RDAG}''; the article states the theorem in its section 6 and prints the proof in its appendix, so the proof is detached from the statement by the article's own arrangement, and both spans are delivered here). No mathematics is written, repaired or re-derived: statement and proof are the article's own text line for line, in the article's own notation ($(G,\mathcal{C})$, $(\mathcal{G},c)$, $c$, $r$, $R$, $m_{v\leftrightarrow w}$, $\ch(v)$, $\mathcal{G}_i$, $\mathcal{G}_{(j\to i)}$, cpeo, (M1), (M2)), and no line of the counted statement or of the counted proof is omitted, substituted or re-flowed. Required context retained uncounted: the article's own definition of $2$-path regularity (lines 695--728, source label \texttt{def:2\_path\_regularity}); its definition of a color elimination-regular (CER) graph and of a symmetric CER graph (lines 743--753); the article's display defining the three conditions (i)--(iii) that the counted statement refers to (lines 2254--2260, source label \texttt{eq:i-iii}); and the four lemmas the counted proof invokes, namely the lemma on equally colored edges (lines 2643--2645, source label \texttt{lem:L\_i}), the lemma identifying the no-induced-$i\to k\leftarrow j$ condition with a perfect elimination ordering (lines 3429--3437, source label \texttt{lem:revtop\_is\_peo}), the lemma on consecutive color blocks in an inverse topological order (lines 3446--3452, source label \texttt{lem:cpeoRDAG}), and the lemma excluding directed edges inside a vertex color class (lines 3479--3482, source label \texttt{lem:no-same-colour-edges}). Editorial formatting only, in addition: an AMS \texttt{amsart} preamble that restates the article's own macro definitions (lines 97--154 of source0.tex, including \texttt{\textbackslash coloredZ}, \texttt{\textbackslash ch}, \texttt{\textbackslash tri}, \texttt{\textbackslash hh}, \texttt{\textbackslash mytransp}, \texttt{\textbackslash tr}, \texttt{\textbackslash BC}, \texttt{\textbackslash cBC}) and the theorem environments the retained lines use; sequential numbering of the retained environments in order of appearance, whereas the article prints the retained items as Theorem 6.1, Definition 6.3, Lemma 6.4, Lemma 8.1, Lemma 8.6, Lemma 8.7 and Lemma 8.8; the article's section-relative equation numbering is not reproduced, so the one displayed relation inside the retained spans is numbered consecutively instead, and every cross-reference inside the retained text resolves to the wrapper's numbering. The article's class file, its package list, its \texttt{.bib} resource and its style files are not shipped; the wrapper is the pinned AMS runtime. No citation occurs inside any retained line, so the wrapper carries no bibliography and no bibliography entry. No asset-inclusion command, no drawing environment and no image asset occurs in this package.

\section*{Retained context: 2-path regularity}

\begin{definition}[$2$-path regularity]\label{def:2_path_regularity}
With the convention above (that is, after relabeling the vertex-color classes so
that the fixed ordering satisfies $\eta_i=i$ for every $i\in[r]$), a colored
graph \((G,\mathcal{C})\) is
\textbf{$2$-path regular with respect to the fixed ordered partition}
$(V_1,\ldots,V_r)$ if the symmetric $2$-path function is constant on
extended-edge color classes, i.e.,
\begin{enumerate}
    \item[(M1)] for all \(\{v, w\}, \{v', w'\} \in \tilde{E}\), we have
    \[
c(v,w)=c(v',w')\quad\implies\quad     m_{v\leftrightarrow w} = m_{v'\leftrightarrow w'}.
    \]
\end{enumerate}
The colored graph \((G,\mathcal{C})\) is
\textbf{symmetric $2$-path regular with respect to the fixed ordered partition}
$(V_1,\ldots,V_r)$ if it is
$2$-path regular with respect to this ordered partition and, additionally, for
any two vertices $v$ and $w$ of the same color, the $2$-path count from $v$ to
$w$ is the same as from $w$ to $v$:
\begin{enumerate}
 \item[(M2)] for all $v,w\in V$, 
 \[
    c(v)=c(w)\quad\implies\quad 
    m_{v\to w}\big|_{F\times F} = m_{w\to v}\big|_{F\times F}.
    \]
\end{enumerate}

We say simply that $(G,\mathcal{C})$ is \textbf{$2$-path regular} if there exists an
ordering of its vertex-color classes with respect to which it is $2$-path
regular. Similarly, $(G,\mathcal{C})$ is \textbf{symmetric $2$-path regular} if there
exists an ordering of its vertex-color classes with respect to which it is
symmetric $2$-path regular. For each such ordering, the vertex-color classes
are relabeled according to the convention above.
\end{definition}


\section*{Retained context: CER and symmetric CER graphs}

\begin{definition}
Let $(G,\mathcal{C})$ be a colored graph with a vertex partition $\mathcal{V}=\{V_1,\ldots,V_r\}$. 

We say that $(G,\mathcal{C})$ is \textbf{color elimination-regular (CER)} if there exists an ordered partition of its vertices, $(V_{\eta_1},\ldots,V_{\eta_r})$, that satisfies two conditions simultaneously:
\begin{itemize}
    \item[(i)] The ordering is a color perfect elimination ordering (cpeo) for $G$. 
    \item[(ii)] The graph $(G,\mathcal{C})$ is $2$-path regular with respect to the same ordering.
\end{itemize}

A \textbf{symmetric CER graph} is a CER graph for which, in addition, $(G,\mathcal{C})$ is symmetric $2$-path regular with respect to the same cpeo.
\end{definition}


\section*{Retained context: the three conditions (i)--(iii)}

\begin{align}\label{eq:i-iii}
\begin{split}
(i)\ & \ \mbox{there is no induced subgraph $i\to k\leftarrow j$ with $i$ and $j$ non-adjacent;}\\
(ii)\ & \ c(i)=c(j)\implies (\mathcal{G}_i,c)\cong (\mathcal{G}_j,c);\\
(iii)\ & \ c(j\to i)=c(l\to k)\implies (\mathcal{G}_{(j\to i)},c)\cong (\mathcal{G}_{(l\to k)},c).
\end{split}
\end{align}


\section*{Retained context: the lemmas used by the proof}

\begin{lemma}\label{lem:L_i}
Assume that $(G,\mathcal{C})$ is a CER graph. If $c(v,w)=c(v',w')=k\in [r+R]\setminus F$, $c(v)\le c(w)$ and $c(v')\le c(w')$, then $c(v)=c(v')$.
\end{lemma}


\begin{lemma}\label{lem:revtop_is_peo}
Let $\mathcal G$ be a DAG with skeleton $G$, and let
$(v_1,\ldots,v_p)$ be an inverse topological ordering,
so that $v_j\to v_i$ implies $i<j$.
Then $\mathcal G$ has no induced subgraph
$i\to k\leftarrow j$ with $i$ and $j$ non-adjacent
if and only if $(v_1,\ldots,v_p)$ is a perfect
elimination ordering of $G$.
\end{lemma}


\begin{lemma}\label{lem:cpeoRDAG}
Let $(\mathcal{G},c)$ be a colored DAG with compatible coloring and assume \eqref{eq:i-iii} (ii).
Then there exists an inverse topological order $\succ$ of $V$ such that each vertex color class forms a consecutive block in $\succ$, i.e.
\[
i\succ j\succ k \ \text{ and }\ c(i)=c(k)\quad\Longrightarrow\quad c(j)=c(i).
\]
\end{lemma}


\begin{lemma}\label{lem:no-same-colour-edges}
Let $(\mathcal{G},c)$ be a colored DAG with compatible coloring and assume \eqref{eq:i-iii} (ii).
Then there are no directed edges between vertices of the same vertex color.
\end{lemma}


\section{The counted unit: Theorem 6.1 and its complete original proof}

\begin{theorem}\label{thm:RDAG}
Let $(G,\mathcal C)$ be a colored graph.
Then $(G,\mathcal C)$ admits an acyclic orientation whose inherited
coloring is compatible and for which (i)--(iii) in \eqref{eq:i-iii}
hold if and only if $(G,\mathcal C)$ is a symmetric CER graph and
$G$ has no edges between vertices of the same color.
\end{theorem}

\begin{proof}[Proof of Theorem~\ref{thm:RDAG}]
\noindent{$(\Rightarrow)$}
Assume that $c$ is a compatible coloring and that $(\mathcal{G},c)$ satisfies \eqref{eq:i-iii}.
By Lemma~\ref{lem:cpeoRDAG} choose an inverse topological order $\succ$ whose color classes are consecutive blocks. 
By \eqref{eq:i-iii} (i) and Lemma~\ref{lem:revtop_is_peo}, the order $\succ$ is a perfect elimination ordering of the skeleton $G$.
By Lemma~\ref{lem:no-same-colour-edges}, $G$ has no edges within a vertex color class; hence the induced ordering of color blocks
is a cpeo.

Relabel the vertex-color classes as $(V_1,\ldots,V_r)$ in the order induced by $\succ$, and let $m$ be the $2$-path function defined with respect to this ordered partition.
For an extended edge $\{v,w\}\in\tilde E$, the $2$-paths
$v-u-w$ involving two ordinary edges and satisfying
$c(u)\le c(v)\wedge c(w)$ correspond exactly to common
children $u$ of $v$ and $w$ (to children of $v$ if $v=w$).
Consequently, for ordinary edge colors $h,h'$ the value $m_{v\leftrightarrow w}(h,h')$ counts common children $u\in\ch(v)\cap\ch(w)$
with unordered pair of incoming edge colors $\{c(v\to u),c(w\to u)\}=\{h,h'\}$.
When $h=h'$, symmetrization counts each child twice;
this common factor does not affect the equality of counts.
The contributions involving loops depend only on the vertex
color when $v=w$, and otherwise on the edge color and the
color of the vertex to which the edge is directed.
They therefore agree on extended-edge color
classes by compatibility, which also determines each child's
vertex color from its incoming edge colors.
Therefore, for vertices $i,j$ of the same color,
$\mathcal{G}_i\cong\mathcal{G}_j$ iff $m_{i\leftrightarrow i}=m_{j\leftrightarrow j}$.
Similarly, for directed edges $j\to i$ and $\ell\to k$ of the same color,
$\mathcal{G}_{(j\to i)}\cong\mathcal{G}_{(\ell\to k)}$ iff $m_{j\leftrightarrow i}=m_{\ell\leftrightarrow k}$.
Hence \eqref{eq:i-iii} (ii)-(iii) are equivalent to (M1), i.e., $2$-path regularity.
Since there are no edges within a vertex color class, (M2) holds trivially, so $(G,\mathcal{C})$ is symmetric CER.

\noindent$(\Leftarrow)$
Assume $(G,\mathcal{C})$ is symmetric CER and that $G$ has no edges between vertices of the same color.
Choose a cpeo from the definition of symmetric CER, relabel the vertex-color classes as $(V_1,\ldots,V_r)$ in this order, and choose any vertex order $\succ$ consistent with it
(so vertices are ordered by color blocks).
Orient every edge from its endpoint in the later color class to its endpoint in the earlier one.
The resulting directed graph $\mathcal G$ is acyclic, since each arrow decreases the color-class index, and $\succ$ is an inverse topological order of $\mathcal G$.
Because there are no edges within any color class, every vertex in a given block has all its later neighbors in later blocks,
and the cpeo condition implies that these later neighbors form a clique; hence $\succ$ is a perfect elimination ordering of $G$.
Applying Lemma~\ref{lem:revtop_is_peo} yields \eqref{eq:i-iii} (i).

Since there are no edges within a vertex-color class, $F=[r]$.
By Lemma~\ref{lem:L_i}, equally colored edges have earlier endpoints of the same vertex color.
The corresponding arrows point towards these endpoints,
so the inherited coloring is compatible.
Moreover, (M1) is exactly constancy of $m_{v\leftrightarrow w}$ on extended-edge color classes, which (by the same common-child
interpretation as above) is equivalent to \eqref{eq:i-iii} (ii)-(iii).
\end{proof}

\end{document}
